\chapter{Our Research Methodology}\label{ch:method}

This chapter describes every approach we built, in the order of the research questions they answer. For each one it
states what is taken from the paper and what is ours. Concrete settings are collected in \cref{ch:setup}, and the code
is described in \cref{ch:impl}.

\section{Research design}\label{sec:design}

\subsection{Two model tracks}

The paper names ``Llama-3.1-8B'' without saying base or instruction-tuned. The two pieces of evidence point in
opposite directions:
\begin{itemize}
\item the repository README links the \emph{base} checkpoint;
\item the repository's only query path is a \emph{chat} call, which a base model cannot serve, because it has no chat
template.
\end{itemize}
We ran both tracks.

\begin{description}[style=nextline,leftmargin=1.2em]
\item[Base track (2026-11-04 onward): faithful to the README.] Llama-3.1-8B base at revision \texttt{d04e592}, queried
with raw completions. The base model was chosen \emph{because the README links it}, before any score was known for it.
An earlier blueprint had chosen Instruct, and we had briefly defended that choice on the grounds that its scores were
nearer the paper's. The user rejected that test on 2026-11-04: choosing a model by how close its score lands to the
target is circular.
\item[Instruct track (2026-11-10 onward): faithful to the repository's query path.] Llama-3.1-8B-Instruct at revision
\texttt{0e9e39f}, queried through the repository's chat call with its system message. It is used in generation,
training (chat template) and evaluation alike. It was adopted for matching the paper after the base track's vanilla
gate failed (\cref{sec:sql-vanilla}), under an approved plan with the rule ``Instruct is treated as the paper's model
for matching; base stays the faithful reproduction''.
\end{description}

\subsection{Arms, seeds and rounds}

\begin{description}[style=nextline,leftmargin=1.2em]
\item[Vanilla.] The untuned model, greedy decoding. It is deterministic, so it is run once.
\item[Table-Specialist (TS).] \Cref{alg:ours-loop}: three rounds with 3,000 new tables per round. Each round's
adapters restart from vanilla and train only on that round's validated pairs. Each round's generators are the previous
round's adapters.
\item[No-validation (NV).] The same loop with the validation step removed: every pair with extractable code is used.
This is the paper's ablation.
\end{description}

Every fine-tuned arm runs with seeds 1, 2 and 3. A seed fixes the table order (each seed takes its own random
permutation of the corpus, and round $r$ uses positions $3000(r-1)$ to $3000r$), the question-sampling RNG, the LLM
request seed, and the training seed. Three seeds was our choice; the paper does not report seeds. Instruct NV was run
for seed~1 only (\cref{sec:sql-instruct}).

\subsection{Principles}

\begin{enumerate}
\item \textbf{Fill gaps with documented choices; do not tune towards the paper's number.} Every unpublished setting is
recorded in \fpath{configs/base.yaml} as ``ours'', with its reason in \fpath{docs/DECISIONS.md}.
\item \textbf{Report both scores.} Repo-faithful scores (the authors' comparison, bugs included) are the headline;
corrected scores are reported beside them.
\item \textbf{Prefer the repository's artefacts.} Prompts, prompt classes, evaluator rules and sampling rules come
from the repository wherever it has them (\cref{sec:provenance}).
\item \textbf{Fix controls and decision rules before seeing results.} Go/no-go thresholds, dev-set selection rules and
same-data controls were written down before the runs they judge.
\end{enumerate}

\section{The Table-Specialist loop as implemented}\label{sec:ours-loop}

\begin{figure}[htbp]
\centering
\begin{adjustbox}{max width=\linewidth}\begin{tikzpicture}[
  box/.style={draw=sea,fill=oursbg,rounded corners=3pt,thick,align=center,font=\sffamily\scriptsize,minimum height=1.05cm,text width=2.5cm},
  repo/.style={box,draw=amber,fill=paperbg},
  arr/.style={-{Stealth[length=4.5pt]},thick,color=slate},
  lab/.style={font=\sffamily\scriptsize,text=slate,fill=white,inner sep=1.5pt}]
  % row 1: questions
  \node[box]  (corpus) at (0,0)    {GitTables corpus\\10,053 tables\\\fpath{ts/corpus/build.py}};
  \node[repo] (qp)     at (3.3,0)  {Question prompt and\\requirement sampler\\(repo, read by AST)};
  \node[box]  (qgen)   at (6.6,0)  {Vanilla writes\\questions, $T=0.7$};
  \node[repo] (cp)     at (9.9,0)  {Few-shot SQL and\\Scala prompts (repo)};
  % row 2: code, validation, training
  \node[box]  (train)  at (0,-2.1)   {LoRA SFT from vanilla:\\SQL and Scala\\adapters};
  \node[box]  (val)    at (6.6,-2.1) {Algorithm 3 with\\our comparator\\\fpath{exec_agree.py}};
  \node[box]  (code)   at (9.9,-2.1) {Previous adapters\\write SQL and Scala;\\greedy, first block};
  % row 3: evaluation and execution
  \node[box]  (eval)   at (0,-4.3)   {Score SQL adapter:\\4 test sets $\times$\\2 prompts};
  \node[repo] (evrepo) at (3.3,-4.3) {NSEvaluator (repo,\\patched) and our\\corrected comparator};
  \node[repo] (sbx)    at (6.6,-4.3) {Spark sandbox (repo;\\open-file patch only)};
  \node[box]  (next)   at (9.9,-4.3) {Adapters become\\the next round's\\generators};
  % shared services band
  \node[draw=slate,fill=rowshade,rounded corners=2pt,align=center,font=\sffamily\scriptsize,text width=12.3cm,minimum height=0.6cm]
    (infra) at (4.95,-5.75) {Shared services: vLLM server $\cdot$ SQLite response cache $\cdot$ GPU locks $\cdot$ Spark validation lock $\cdot$ resume from files};
  \draw[arr] (corpus) -- (qp); \draw[arr] (qp) -- (qgen); \draw[arr] (qgen) -- (cp);
  \draw[arr] (cp) -- (code); \draw[arr] (code) -- (val);
  \draw[arr] (val) -- node[lab,above=1pt]{kept pairs (TS) / all pairs (NV)} (train);
  \draw[arr] (sbx) -- (val);
  \draw[arr] (train) -- (eval); \draw[arr] (evrepo) -- (eval);
  \draw[arr] ($(train.south)+(0.75,0)$) |- ($(next.north)+(-0.6,0.42)$) -- ($(next.north)+(-0.6,0)$);
  \draw[arr,dashed] (next.east) -- ++(0.35,0) |- (code.east);
\end{tikzpicture}
\end{adjustbox}
\caption{Our implementation of one round (\fpath{ts/loop.py}). Teal boxes are our code, amber boxes are repository
artefacts used unchanged. The response cache and the GPU and CPU locks make every stage resumable and keep one job per
GPU.}
\label{fig:our-pipeline}
\end{figure}

\begin{algorithm}[htbp]
\caption{One round of our TS / NV loop for NL-to-SQL (\fpath{ts/loop.py: run_round})}\label{alg:ours-loop}
\KwIn{seed $s$; round $r$; adapters $A^{(r-1)}_{\text{SQL}}, A^{(r-1)}_{\text{Scala}}$ (none in round 1); arm
$\in\{\text{TS},\text{NV}\}$}
$\mathcal{T}_r \leftarrow \operatorname{perm}_s(\mathcal{T})[3000(r-1):3000r]$\;
\ForEach{$t\in\mathcal{T}_r$}{
  draw requirements with RNG$(1000s+r)$ (repo sampler; redraw if empty)\;
  $x \leftarrow$ vanilla$(\text{question prompt}(t))$ at $T=0.7$, request seed $s$\Comment*[r]{ours: vanilla writes questions}
}
$y \leftarrow A^{(r-1)}_{\text{SQL}}(\text{few-shot SQL prompt})$;\ $y' \leftarrow A^{(r-1)}_{\text{Scala}}(\text{few-shot
Scala prompt})$, greedy\;
keep the first complete code block of each reply\Comment*[r]{ours}
\eIf{arm $=$ TS}{
  $\mathcal{D}_r \leftarrow \{(x,y,y') : V_{\text{ours}}(t,y,y')=1\}$\Comment*[r]{\cref{alg:ours-validate}}
}{
  $\mathcal{D}_r \leftarrow \{(x,y,y') : y \neq \varnothing,\ y'\neq\varnothing\}$\;
}
build examples in the repo's NL-to-SQL and NL-to-Scala training formats\;
train $A^{(r)}_{\text{SQL}}$ and $A^{(r)}_{\text{Scala}}$ \emph{from vanilla} with LoRA\Comment*[r]{restart, paper Sec.~6.5}
evaluate $A^{(r)}_{\text{SQL}}$ on 4 test sets $\times$ 2 prompts\;
\end{algorithm}

\subsection{Questions and programs}

\begin{itemize}
\item \textbf{Questions} use the repository's question prompt and requirement sampler, read from
\fpath{table2question_generationv_v4.py} by parsing the script's syntax tree, without running it. The vanilla model
writes the questions in every round, at temperature 0.7, with a token limit of 128 (base) or 256 (Instruct). The paper
does not say which model writes questions; the repository's question files are named for Llama, in a fresh table slice
per round.
\item \textbf{Programs} use the repository's few-shot generation prompts for SQL and Spark-Scala, read the same way.
Decoding is greedy (temperature 0).
\begin{itemize}
\item Base: limits of 128 tokens (SQL) and 256 (Scala).
\item Instruct: limits of 256 and 512, because chat replies add prose around the code.
\item We keep the \emph{first} complete fenced code block. A base model keeps writing after its answer, and the
repository's parser takes the last block. A block cut off by the token limit has no closing fence and counts as no
code.
\end{itemize}
\item \textbf{Adapter names.} Each adapter is served by vLLM~\cite{kwon2023vllm} under a name derived from a hash of
its weights. The response cache keys on the model name, so fixed names such as \texttt{sql} would have let one arm
silently reuse another arm's cached answers.
\end{itemize}

\subsection{Validation: our comparator}\label{sec:comparator}

The repository's validation scripts run both languages but never compare their results, so the comparison is ours.
Let a result be a list of rows. For a cell $v$, define the text form $\tau(v)$ as the trimmed string, with
$\tau(\text{NULL})=\tau(\text{NaN})=\text{``''}$, and map the strings \texttt{true}/\texttt{false} to
\texttt{1}/\texttt{0}. Two cells agree if
\begin{equation}
  a \simeq b \iff \tau(a)=\tau(b)\ \ \lor\ \ \bigl|\,\operatorname{float}(\tau(a))-\operatorname{float}(\tau(b))\,\bigr|
  \le \epsilon,\qquad \epsilon=10^{-6}.
\end{equation}
Two rows agree if they have equal length and agree cell by cell, in column order. Let $E(R)$ be result $R$ with
all-empty rows removed. Two results agree if
\begin{equation}
  R \equiv R' \iff |E(R)|=|E(R')| \ \land\ \exists\ \text{a bijection } \sigma: E(R)\to E(R') \text{ with } \rho \simeq
  \sigma(\rho)\ \ \forall \rho,
\end{equation}
found greedily. Row order and column names are ignored, because Spark gives no order without a sort and the two
languages name computed columns differently. A pair whose SQL returns no non-empty row on \emph{all} 10 samples is
rejected as \texttt{always\_empty}: it has shown no agreement about any data. A literal reading of
\cref{alg:paper-exec} would keep it.

Five rules were added after a hand inspection of pairs (\cref{sec:fail-validator}). Each fixed a case where identical
logic was being rejected:
\begin{enumerate}
\item SQL samples are written to CSV and read back with type inference before loading SQLite, as the repository's
\fpath{execute.py} does. Our first version kept all text, so \texttt{MAX("Weight")} compared strings.
\item All-empty rows are ignored. Spark writes no line for a one-column null row.
\item Booleans compare equal to 1 and 0.
\item Spark's \texttt{\textbackslash"} escaping is read correctly.
\item Adapter names are derived from weight hashes (above).
\end{enumerate}
SQL runs first, because it takes milliseconds in SQLite. Scala stops at the first disagreeing sample, because every
sample must agree. Neither shortcut can change a verdict.

\begin{algorithm}[htbp]
\caption{Our validation of one pair (\fpath{ts/validate/exec_agree.py: validate_pair})}\label{alg:ours-validate}
\KwIn{table $t$ (read as text); SQL $y$; Scala $y'$}
$S_0,\dots,S_9 \leftarrow$ \texttt{t.sample(frac=0.5, random\_state=}$k$\texttt{)}\;
\For{$k\leftarrow 0$ \KwTo $9$}{$R_k \leftarrow \operatorname{SQLite}(y, \operatorname{reinfer}(S_k))$; \lIf{error}{\Return{sql\_failed}}}
\lIf{$E(R_k)=\varnothing$ for all $k$}{\Return{always\_empty}}
\For{$k\leftarrow 0$ \KwTo $9$}{
  $R'_k \leftarrow \operatorname{Spark}(y', S_k)$ in the Docker sandbox\;
  \lIf{failure or unreadable output}{\Return{scala\_failed}}
  \lIf{$R_k \not\equiv R'_k$}{\Return{disagree}}
}
\Return{agree}
\end{algorithm}

\subsection{Training}

Both adapters are trained with LoRA (\cref{eq:lora}) and completion-only loss (\cref{eq:sft}), using TRL's SFT trainer
with PEFT, from the vanilla weights each round:
\begin{itemize}
\item rank $r=16$, $\alpha=32$, dropout 0.05, all linear layers;
\item learning rate $2\times10^{-4}$, cosine schedule, 3 epochs, bf16;
\item maximum length 4,096 tokens.
\end{itemize}
The paper's batch of 1\% of the data is built by gradient accumulation of single examples: accumulation $=\lceil
0.01\,|\mathcal{D}_r|\rceil$. That gives about 300 optimiser steps for any data size (3 epochs of about 100 steps each).
Examples use the repository's prompt classes for NL-to-SQL and NL-to-Scala, with the zero-shot prompt. The repository's
own few-shot training defaults raise an error, so its training format cannot be reproduced exactly
(\cref{sec:fail-repo}). On Instruct, examples are in TRL's conversational prompt--completion form with the repository's
system message, so training sees the same template as serving.

Checkpoints are saved about every twelfth of a run, or every 5 minutes for large adapters, with the trainer keeping
only the latest. A crash therefore loses at most about 5 minutes of training. Checkpoints were copied aside before
every resume (\cref{app:incidents}).

\section{Evaluation and statistics}\label{sec:scoring}

\subsection{Test sets and prompts}

The SQL adapter of every round is scored on four rebuilt test sets (\cref{sec:testsets}) with two prompts:
\begin{itemize}
\item \textbf{zero-shot}: the repository's default test prompt, from its prompt class;
\item \textbf{few-shot}: the repository's four-example generation prompt.
\end{itemize}
The paper does not state its test prompt, so the comparison with the paper uses the zero-shot rows. This choice
followed the repository's default; it was \emph{not} made because zero-shot scores were closer to the paper's. On base,
few-shot was closer on 3 of 4 sets.

\subsection{Repo-faithful and corrected scores}

The repo-faithful score runs the repository's \texttt{NSEvaluator} compiled from source with checked patches
(\cref{sec:provenance}). Its list comparison \texttt{are\_lists\_equal} has three defects, which the corrected score
(\fpath{ts/eval/compare.py}) fixes, keeping its numeric tolerance unchanged:
\begin{enumerate}
\item It returns \emph{equal} at the first cell that differs but is numerically close, skipping every later cell.
\item A differing cell whose gold value is 0 raises a division error, which counts as a mismatch.
\item Rows are compared in order even when the gold query has no \texttt{ORDER BY}, where row order is undefined.
\end{enumerate}
The corrected comparison treats two cells as equal if they are identical, or if both are numbers with absolute
difference below 0.01 and relative difference below 0.01 (absolute difference alone when the gold value is 0). Results
are compared as multisets unless the gold SQL contains \texttt{ORDER BY}. The corrected score also does not lowercase
gold SQL, because lowercasing breaks non-ASCII column names such as \texttt{"Álbum"}.

\subsection{Statistics}\label{sec:stats}

Each cell reports a 95\% Wilson interval~\cite{wilson1927} for its accuracy $\hat p$ over $n$ cases:
\begin{equation}
  \frac{\hat p + \frac{z^2}{2n} \pm z\sqrt{\frac{\hat p(1-\hat p)}{n}+\frac{z^2}{4n^2}}}{1+\frac{z^2}{n}},\qquad z=1.96.
\end{equation}

Comparisons between models use a paired bootstrap over test cases~\cite{efron1993bootstrap}
(\fpath{ts/analysis/significance.py}).
\begin{itemize}
\item Let $a_i$ and $b_i$ be the correctness of models A and B on case $i$, each averaged over the seeds both models
have.
\item For cells $c=1..C$ ($C=8$, or 4 for the zero-shot mean) with case sets $I_c$, the statistic is
\begin{equation}
  \Delta \;=\; \frac{1}{C}\sum_{c=1}^{C}\frac{1}{|I_c|}\sum_{i\in I_c}(a_i-b_i).
\end{equation}
\item Cases are resampled with replacement within each cell, $B=10{,}000$ times with a fixed RNG seed.
\item The 95\% interval is the 2.5th and 97.5th percentiles of $\Delta^\ast$. The two-sided $p$-value is $2\min\bigl(\Pr[\Delta^\ast\le 0],\Pr[\Delta^\ast\ge
0]\bigr)$, capped at 1.
\end{itemize}
The case bootstrap does not cover seed-to-seed variation, so per-seed differences are always reported beside it. With
three seeds, a finding is only called consistent when all three seeds agree in sign.

\section{Extension 2: learning from rejected programs}\label{sec:method-ext2}

\begin{oursbox}[title={Our extension (not in the paper)}]
\textbf{Motivation.} TS keeps only about 15\% of pairs, and a hand check suggested that about one rejected pair in three
holds usable SQL (\cref{sec:fail-validator}). Two questions follow. Does sampling more candidates per question recover
validated data (\emph{bestof4})? Can the rejected programs themselves teach the model what not to write (\emph{DPO})?
\end{oursbox}

\textbf{Labelling (base, round~1, same questions and tables as TS round~1).}
\begin{itemize}
\item Each question has 4 candidate pairs: TS's greedy pair (its verdict reused) plus 3 pairs sampled from vanilla at
$T=0.7$, with request seeds $100s+j$.
\item The \emph{chosen} program is the first candidate, in order, whose SQL and Scala pass Algorithm~3 unchanged.
\item A candidate is \emph{rejected} if its SQL rows differ from the chosen SQL's on some sample, or it fails, or it is
always empty. It is never judged by its own Scala, so that a correct SQL whose Scala failed does not become a negative.
\item Spark runs only until a question has a chosen program; later candidates are labelled by SQL results alone.
\end{itemize}

\textbf{Models.}
\begin{itemize}
\item \emph{bestof4}: SFT from vanilla on one chosen program per question, with the TS settings.
\item \emph{DPO}~\cite{rafailov2023dpo}: bestof4, then preference training on (chosen, rejected) pairs with an added
SFT term on the chosen program (the RPO loss~\cite{pang2024rpo}):
\end{itemize}
\begin{equation}
  \mathcal{L} \;=\; -\log\sigma\!\Bigl(\beta\Bigl[\log\frac{\pi_\theta(y_w\mid x)}{\pi_{\text{ref}}(y_w\mid x)} -
  \log\frac{\pi_\theta(y_l\mid x)}{\pi_{\text{ref}}(y_l\mid x)}\Bigr]\Bigr) \;-\; \lambda\,\frac{\log\pi_\theta(y_w\mid
  x)}{|y_w|},
\end{equation}
with $\pi_{\text{ref}}$ the frozen bestof4 adapter, $\beta=0.1$, $\lambda=1$, LoRA learning rate $2\times10^{-5}$,
cosine, 10\% warm-up, 2 epochs and 16 pairs per step. Five gentler variants were tuned on dev sets only, under a
selection rule fixed in the configuration before any dev result (\cref{sec:ext2}).

\section{Extension 3: repair-TS}\label{sec:method-repair}

\begin{oursbox}[title={Our extension (not in the paper)}]
\textbf{Idea.} Give each rejected pair one repair attempt, guided only by what its own program did on the table, and
then validate it again with Algorithm~3 \emph{unchanged}. A program never sees the other language's code or output, so
agreement after repair is still two independent programs agreeing. The validator's standard is not lowered.
\end{oursbox}

\begin{figure}[htbp]
\centering
\begin{adjustbox}{max width=\linewidth}\begin{tikzpicture}[
  node distance=0.5cm and 0.6cm,
  box/.style={draw=sea,fill=oursbg,rounded corners=3pt,thick,align=center,font=\sffamily\scriptsize,minimum height=0.9cm,text width=2.5cm},
  rep/.style={box,draw=coral,fill=coral!8},
  dec/.style={diamond,draw=amber,fill=paperbg,thick,aspect=2,inner sep=1pt,align=center,font=\sffamily\scriptsize,text width=1.6cm},
  arr/.style={-{Stealth[length=4.5pt]},thick,color=slate}]
  \node[box,text width=3.1cm] (rej) {Rejected pair\\$(x,\,y_{\text{SQL}},\,y_{\text{Scala}})$\\reason, sample $k$};
  \node[rep,above right=0.15cm and 0.8cm of rej] (sql) {Repair SQL\\\scriptsize sql\_failed, always\_empty};
  \node[rep,right=0.8cm of rej] (both) {Repair both\\\scriptsize disagree};
  \node[rep,below right=0.15cm and 0.8cm of rej] (sca) {Repair Scala\\\scriptsize scala\_failed};
  \node[box,right=0.8cm of both,text width=3.0cm] (fb) {Feedback from the program's own run on $S_k$:\\error $|$ ``no rows'' $|$ its first 5 rows};
  \node[box,right=0.6cm of fb,text width=2.4cm] (gen) {Instruct repairer\\greedy, then $\le$3 samples at $T{=}0.7$};
  \node[dec,below=0.9cm of gen] (v) {Algorithm 3\\unchanged};
  \node[box,left=1.0cm of v] (rec) {Recovered pair\\joins training data};
  \node[box,right=0.6cm of v,text width=1.6cm,draw=slate,fill=rowshade] (drop) {Discarded};
  \draw[arr] (rej) -- (sql); \draw[arr] (rej) -- (both); \draw[arr] (rej) -- (sca);
  \draw[arr] (sql) -| (fb); \draw[arr] (both) -- (fb); \draw[arr] (sca) -| (fb);
  \draw[arr] (fb) -- (gen); \draw[arr] (gen) -- (v);
  \draw[arr] (v) -- node[above,font=\sffamily\scriptsize]{agree} (rec);
  \draw[arr] (v) -- node[above,font=\sffamily\scriptsize]{no} (drop);
\end{tikzpicture}
\end{adjustbox}
\caption{Repair-TS. The failure type decides which side is repaired; each side's feedback comes only from its own
execution on the sample where the pair was rejected. Recovered pairs join the kept pairs for training.}
\label{fig:repair-flow}
\end{figure}

\paragraph{Which side is repaired.} The failure type decides:
\begin{itemize}
\item \texttt{sql\_failed} and \texttt{always\_empty}: repair the SQL.
\item \texttt{scala\_failed}: repair the Scala.
\item \texttt{disagree}: repair both separately, then try (new SQL, old Scala), (old SQL, new Scala) and (new, new) in
that order, keeping the first that passes.
\item \texttt{no\_code}: not repaired.
\end{itemize}

\paragraph{Feedback.}
\begin{itemize}
\item The SQLite or Spark error message (at most 3 lines, file paths removed).
\item Or ``returns no rows''.
\item Or the program's own first 5 result rows on the sample where the pair was rejected.
\end{itemize}

\paragraph{Prompts.}
\begin{itemize}
\item \emph{Base repairer}: the repository's generation prompt, then the failed code block, its feedback, and the
repository's own correction phrase from \fpath{nss_task_descriptor.py}: ``Upon review, code is incorrect and should be
rewritten as follows:''.
\item \emph{Instruct repairer}: the same content in one user message, asking for the corrected code only.
\end{itemize}

\paragraph{Retries.} Repair is greedy first. If the reply is empty or identical to the failed code, up to 3 samples at
$T=0.7$ are tried (request seeds $100s+1..3$).

\paragraph{Go / no-go rule.} Fixed before the pilot: kept pairs must rise by at least 30\%, and at least 80\% of
hand-checked recovered pairs must be correct. The base repairer failed the rule. Instruct passed it and is the
repairer in every full run (\cref{sec:fail-baserepair}). On base, repair-TS therefore uses a second model for repairs,
which matters for interpretation (\cref{ch:discussion}).

\paragraph{Rounds.} Round~1 repairs TS round~1's rejected pairs. Rounds 2--3 follow \cref{alg:ours-loop} with repair-TS's
own previous adapters as generators, then repair. Both adapters train from vanilla on kept plus recovered pairs.

\begin{algorithm}[htbp]
\caption{Repair-TS for one round (\fpath{ts/repair.py}, \fpath{ts/repair_loop.py})}\label{alg:repair}
\KwIn{validated pairs of the round; repairer $M_R$ (Instruct)}
\ForEach{rejected pair $(x,y,y')$ with reason $\in$\{sql\_failed, always\_empty, scala\_failed, disagree\}}{
  $k \leftarrow$ index of the first rejecting sample\;
  feedback for each side to repair $\leftarrow$ its own error, ``no rows'', or its first 5 rows on $S_k$\;
  new code $\leftarrow M_R$(repo prompt $\Vert$ failed code $\Vert$ feedback), greedy, then up to 3 samples if unchanged\;
  \ForEach{candidate (SQL, Scala) in the fixed order}{
    \lIf{\cref{alg:ours-validate} returns agree}{recover the candidate; \textbf{break}}
  }
}
train SQL and Scala adapters from vanilla on kept $\cup$ recovered\;
\end{algorithm}

\section{Controls for repair-TS}\label{sec:method-controls}

\paragraph{Same-data control} (\fpath{ts/data_control.py}). For seed $s$, train on TS round~1's kept pairs plus $N$
extra validated pairs, where $N$ is repair-TS round~1's recovered count for that seed. The extra pairs are drawn with
\texttt{Random(}$s$\texttt{)} from TS round~1 of donor seed $s \bmod 3 + 1$, excluding tables seed $s$ used. They come
from the same generator and passed the same check. The control therefore differs from repair-TS only in \emph{where}
the extra pairs come from, and from TS only in \emph{how many} pairs there are.

\paragraph{bestof4 as a sampling control.} bestof4 recovers data by sampling rather than by feedback-guided repair, on
the same questions.

\section{Self-consistency at test time}\label{sec:method-sc}

As an established test-time method~\cite{wang2022selfconsistency,wang2018egd}, each test case gets the greedy answer
plus $m=8$ samples at temperature $T$. Let $\mathcal{A}$ be the set of answers that execute and return a non-empty
result, and group answers by identical execution result. The vote returns the first answer of the largest group,
\begin{equation}
  \hat y \;=\; \arg\max_{g}\ |g|,\qquad g \subseteq \mathcal{A},\ \ \text{ties broken by order (greedy first)},
\end{equation}
or the greedy answer if $\mathcal{A}=\varnothing$. The gold answer is never used. $T$ was chosen on dev sets
(\cref{sec:ext-sc}). Self-consistency is reported \emph{beside} greedy results, never in their place: the paper's
numbers are greedy, and the vote costs 9 generations per case.

\section{The paper's other tasks}\label{sec:method-other}

\subsection{NL-to-Scala}

Every TS and repair-TS round already trained a Scala adapter on the validated pairs, so NL-to-Scala needs only scoring.
19 Instruct models are scored: vanilla, plus TS and repair-TS for 3 seeds $\times$ 3 rounds. Scoring follows the
repository's \fpath{nss_evaluator.py} (\fpath{ts/eval/scala_eval.py}):
\begin{itemize}
\item \textbf{Prompts}: zero-shot from the NSS prompt class, few-shot from \fpath{nl2sparkscala_generation_v6.py}.
\item \textbf{Code}: the \emph{last} \texttt{scala} block. It must read \texttt{data.csv} and write
\texttt{.write.option("header","true").csv("result")}, or it is not run.
\item \textbf{Execution}: Spark runs on the full table. BIRD tables are exported from the case database.
\item \textbf{Repo-faithful verdict}: exact, order-sensitive list equality after \texttt{to\_numeric} and
\texttt{fillna("None")}.
\item \textbf{Corrected verdict}: the corrected comparison on the same results.
\end{itemize}
Identical (code, table) pairs run once, through a shared execution cache.

\subsection{Error detection}\label{sec:method-ed}

\begin{oursbox}[title={Choices the paper leaves open}]
The paper and repository give the prompts and Algorithm~2, but not which model writes the errors, nor how rounds are
chained for this task. Our choices mirror the SQL study:
\begin{itemize}
\item \emph{vanilla} writes the errors, at $T=0.7$;
\item the \emph{validator} is the previous round's error-detection adapter (vanilla in round~1), greedy;
\item each round trains from vanilla on its own kept columns.
\end{itemize}
\end{oursbox}

\paragraph{Per round, per seed} (\fpath{ts/ed_loop.py}):
\begin{enumerate}
\item On the round's 3,000 tables, one random column per table receives a typo, from the repository's
\fpath{table2typo_genneration_v1.py}.
\item Another random column receives a semantic error, from \fpath{table2semerr_genneration_v6.py}. Columns with a
value over 50 characters are skipped, as in the repository.
\item An error that is empty, or already present in the column, is dropped.
\item Clean columns come from the first 1,000 tables, at most 2 per table, as in the repository's negative script.
\item \Cref{alg:paper-perm} validates each candidate with 5 shuffles.
\item Kept positives are labelled with the error; kept clean columns with NULL.
\end{enumerate}
Examples use the repository's test format (\texttt{EDPromptGenerator}, zero-shot). Only columns whose prompt fits the
length limit are used.

Scoring follows \fpath{ed_evaluator.py} (\cref{eq:edf1}). We also report AutoTest's own protocol, in which predictions
of debatable values count neither way. The test columns are those whose repository prompt fits 70\% of 4,096 tokens.
That rule gives exactly the paper's 1,126 (ST) and 1,081 (RT) columns.

\paragraph{Repair-ED and its controls} (\fpath{ts/repair_ed.py}). Each error the validator rejects gets one more
attempt from vanilla.
\begin{itemize}
\item \textbf{Repair prompt.} The repository's generation prompt, the model's own answer, then what the validator did
on the 5 shuffles: how many found the error, how many found nothing, and up to 3 other values it named, each cut to 60
characters. It ends with our phrase, modelled on the repository's: ``Upon review, the error is not obvious enough and
should be rewritten as follows:''.
\item \textbf{Validation.} Algorithm~2 runs again, unchanged, with the same validator.
\item \textbf{Reuse.} Repair-ED reuses TS's generated errors, and in round~1 TS's validation, both of which are
identical by construction.
\item \textbf{Best-of-2 control.} The same prompt, sampled again with no feedback (request seed $+1000$).
\item \textbf{Same-data control.} TS round~1's kept columns, plus as many extra kept errors of each kind as repair-ED
recovered, taken from a donor seed.
\end{itemize}

\subsection{NL-to-R}\label{sec:method-r}

The repository has an NL-to-R generation prompt and evaluator, but no validation script. Our design
(\fpath{ts/r_loop.py}, \fpath{ts/validate/r_agree.py}, \fpath{ts/eval/r_eval.py}) pairs R with SQL:
\begin{itemize}
\item \textbf{Validation.} Algorithm~3 with R in place of Scala: same samples, same comparator. R's unquoted
\texttt{NA} is read as an empty cell.
\item \textbf{Data.} Questions are TS's own, and round~1 SQL is TS round~1's. Both are identical by construction.
\item \textbf{Training.} Both SQL and R adapters train each round.
\item \textbf{Scoring.} The repository's \texttt{nr\_evaluator} rules: the \emph{first} \texttt{r} block, no format
check, exact list equality.
\end{itemize}
It runs in an R sandbox we built from the repository's recipe plus \texttt{dplyr}. It was smoke-tested on GPU before
the full runs; results are in \cref{sec:results-r}.
